\documentclass[11pt]{article}
\usepackage[a4paper,margin=1in]{geometry}
\usepackage{fontspec}
\usepackage[boldfont=false]{xeCJK}
\setCJKmainfont{FandolSong-Regular.otf}
\usepackage{amsmath}
\usepackage{amssymb}
\usepackage{array}
\usepackage{booktabs}
\usepackage{graphicx}
\usepackage{adjustbox}
\usepackage{enumitem}
\setlist[itemize]{topsep=2pt,partopsep=0pt,parsep=0pt,itemsep=2pt,leftmargin=1.4em}
\setlist[enumerate]{topsep=2pt,partopsep=0pt,parsep=0pt,itemsep=2pt,leftmargin=1.6em}
\usepackage{parskip}
\usepackage{fvextra}
\fvset{breaklines=true,breakanywhere=true,fontsize=\small}
\setlength{\parindent}{0pt}

\begin{document}

\begin{center}
  {\LARGE\bfseries Carboxylic acids and derivatives}\\[0.35em]
  {\large A Level Chemistry}
\end{center}
\vspace{0.6em}

\section*{Carboxylic acids}

A \textbf{carboxylic acid} 羧酸 has the $\text{–COOH}$ (carboxyl) functional group. It is a weak acid.

\subsection*{Making carboxylic acids}

\begin{itemize}
\item \textbf{oxidation} 氧化 of a primary \textbf{alcohol} 醇 or an \textbf{aldehyde} 醛 with acidified $\text{K}_2\text{Cr}_2\text{O}_7$ or $\text{KMnO}_4$, with \textbf{reflux} 回流 (so it is fully oxidised).

\item \textbf{hydrolysis} 水解 of a \textbf{nitrile} 腈 with dilute acid or alkali, then acidifying.

\item hydrolysis of an \textbf{ester} 酯 with dilute acid or alkali and heat, then acidifying.

\end{itemize}

\subsection*{Reactions of carboxylic acids}

These reactions all show that carboxylic acids are acids:

\begin{itemize}
\item \textbf{with a reactive metal} (a \textbf{redox} 氧化还原 reaction): gives a \textbf{salt} 盐 and hydrogen.

\end{itemize}

\[
2\text{CH}_3\text{COOH} + \text{Mg} \rightarrow (\text{CH}_3\text{COO})_2\text{Mg} + \text{H}_2
\]

\begin{itemize}
\item \textbf{with an alkali} (a \textbf{neutralisation} 中和): gives a salt and water.

\end{itemize}

\[
\text{CH}_3\text{COOH} + \text{NaOH} \rightarrow \text{CH}_3\text{COONa} + \text{H}_2\text{O}
\]

\begin{itemize}
\item \textbf{with a carbonate} 碳酸盐: gives a salt, water and carbon dioxide. The fizzing of $\text{CO}_2$ is a test for a carboxylic acid.

\end{itemize}

\[
2\text{CH}_3\text{COOH} + \text{Na}_2\text{CO}_3 \rightarrow 2\text{CH}_3\text{COONa} + \text{H}_2\text{O} + \text{CO}_2
\]

Two reactions change the functional group:

\begin{itemize}
\item \textbf{esterification} 酯化 with an alcohol (concentrated $\text{H}_2\text{SO}_4$ catalyst) gives an ester.

\item \textbf{reduction} 还原 by $\text{LiAlH}_4$ gives a primary alcohol.

\end{itemize}

\section*{Esters}

An ester has the $\text{–COO–}$ group. It often smells sweet or fruity.

\subsection*{Making esters}

An ester forms in a \textbf{condensation} 缩合 reaction between an alcohol and a carboxylic acid, with concentrated $\text{H}_2\text{SO}_4$ as catalyst. A water molecule is lost, and the reaction is reversible:

\[
\text{CH}_3\text{COOH} + \text{C}_2\text{H}_5\text{OH} \rightleftharpoons \text{CH}_3\text{COOC}_2\text{H}_5 + \text{H}_2\text{O}
\]

\subsection*{Hydrolysis of esters}

Hydrolysis splits the ester back apart. The conditions change the products:

\begin{itemize}
\item \textbf{dilute acid} and heat: reversible. Gives back the carboxylic acid and the alcohol.

\item \textbf{dilute alkali} and heat: not reversible. Gives the alcohol and the \textbf{salt} of the carboxylic acid (the carboxylate ion).

\end{itemize}


\end{document}
